\chapter{Espaces vectoriels}\label{ch:b1:vspaces}

L'algèbre linéaire commence ici~: les axiomes des \omterm{def:b1:vspaces:def}{espaces vectoriels}
isolent ce que $\R^2$, $\R^3$, les espaces de \omterm{def:b1:poly:def}{polynômes} et les espaces
de fonctions ont en commun --- on peut additionner, et multiplier par
un scalaire. Deux chapitres construisent la théorie
(le \cref{ch:b1:findim} y ajoute la dimension)~; le langage qu'ils
installent --- \omterm{def:b1:vspaces:span}{sous-espace engendré}, \omterm{def:b1:vspaces:free}{famille libre}, \omterm{def:b1:vspaces:free}{base}, \omterm{def:b1:vspaces:sum}{somme
directe} --- est le pain quotidien de tous les chapitres qui suivent.
Dans tout ce qui suit, $K$ désigne $\R$ ou $\C$ (les
\emph{scalaires}).

\section{Définition et exemples}

\begin{definition}[Espace vectoriel]\label{def:b1:vspaces:def}
Un \emph{$K$-espace vectoriel}\index{espace vectoriel} est un \omterm{def:b1:logic:sets}{ensemble}
$E$ muni d'une addition faisant de $(E, +)$ un \omterm{def:b1:structures:group}{groupe abélien} (de
neutre noté $0_E$ ou $0$), et d'une multiplication externe $K \times E
\to E$ telle que, pour tous $\lambda, \mu \in K$ et $x, y \in E$~:
\[
\lambda(x + y) = \lambda x + \lambda y,\quad
(\lambda + \mu) x = \lambda x + \mu x,\quad
\lambda(\mu x) = (\lambda\mu) x,\quad
1\,x = x .
\]
Conséquences~: $0\,x = 0_E$, $\lambda\,0_E = 0_E$, $(-1)x = -x$, et
$\lambda x = 0_E \implies \lambda = 0$ ou $x = 0_E$ (multiplier par
$\lambda^{-1}$).
\end{definition}

\begin{proof}[Démonstration des conséquences]
Pour $0\,x = 0_E$~: de $(0 + 0)x = 0x + 0x$ et $(0+0)x = 0x$, on
simplifie $0x$ dans le \omterm{def:b1:structures:group}{groupe} $(E, +)$. Pour $\lambda\,0_E$~: même
astuce sur $\lambda(0_E + 0_E)$. Pour $(-1)x$~: ajoutons $x$,
\[
x + (-1)x = 1\,x + (-1)x = \bigl(1 + (-1)\bigr)x = 0\,x = 0_E ,
\]
donc $(-1)x$ est l'opposé de $x$. Enfin si $\lambda x = 0_E$ avec
$\lambda \neq 0$~: multiplions par $\lambda^{-1}$ (les scalaires
forment un \omterm{def:b1:structures:field}{corps}) et utilisons les deux axiomes
$\lambda^{-1}(\lambda x) = (\lambda^{-1}\lambda)x = 1x = x$
ainsi que $\lambda^{-1}0_E = 0_E$~: $x = 0_E$. Aussi modestes
soient-elles, ces quatre règles servent en silence à chaque page qui
suit --- et la dernière est exactement l'endroit où l'on a besoin d'un
\omterm{def:b1:structures:field}{corps}~: sur les scalaires $\Z$, l'«~espace~» $\Z/2\Z$ la mettrait en
défaut avec $2\,x = 0$.
\end{proof}

\begin{example}\label{ex:b1:vspaces:examples}
$K^n$ (opérations coordonnée par coordonnée)~; les \omterm{def:b1:poly:def}{polynômes} $K[X]$~;
les fonctions $\mathcal{F}(A, K)$ d'un \omterm{def:b1:logic:sets}{ensemble} quelconque $A$ dans
$K$ (opérations point par point) --- qui contient les fonctions
\omterm{def:b1:continuity:continuous}{continues}, les suites $\mathcal{F}(\N, \R)$, etc.~; $\C$ comme
$\R$-espace vectoriel. Dans chaque cas, les axiomes sont hérités de
ceux de $K$.
\end{example}

\begin{definition}[Sous-espace vectoriel]\label{def:b1:vspaces:subspace}
$F \subseteq E$ est un \emph{sous-espace
vectoriel}\index{sous-espace vectoriel} lorsque $0_E \in F$ et que $F$
est stable par addition et par multiplication externe --- de façon
équivalente~:
\[
F \neq \emptyset
\qquad\text{et}\qquad
\forall x, y \in F,\ \forall \lambda \in K,\quad
x + \lambda y \in F .
\]
Un sous-espace vectoriel est lui-même un \omterm{def:b1:vspaces:def}{espace vectoriel}. Toute
intersection de sous-espaces est un sous-espace~; une réunion ne l'est
presque jamais (même démonstration qu'à
\cref{exo:b1:structures:6}).
\end{definition}

\begin{example}\label{ex:b1:vspaces:subspaces}
Dans $\mathcal{F}(\R, \R)$~: les fonctions \omterm{def:b1:continuity:continuous}{continues}, les fonctions
\omterm{def:b1:derivative:def}{dérivables}, les \omterm{def:b1:poly:def}{polynômes} de degré $\leq n$ (notés $K_n[X]$ à
l'\omterm{def:b1:topology:closure}{intérieur} de $K[X]$), les solutions d'une équation différentielle
linéaire homogène (le \cref{thm:b1:diffeq:homogeneous2} disait
exactement cela). Contre-exemples~: $\{f : f(0) = 1\}$ (pas de
vecteur nul), les \omterm{def:b1:poly:def}{polynômes} de degré exactement $n$ (non stable par
addition).
\end{example}

\begin{example}[Sous-espace ou non~: quatre verdicts, argumentés]
\label{ex:b1:vspaces:subspaceornot}
Dans l'espace des suites réelles~:
\begin{itemize}
  \item $\{u : u \text{ bornée}\}$ \emph{est} un \omterm{def:b1:vspaces:subspace}{sous-espace}~: $0$ est
        bornée, et si $\abs{u_n} \leq M$, $\abs{v_n} \leq M'$,
        alors $\abs{u_n + \lambda v_n} \leq M + \abs\lambda M'$.
  \item $\{u : u_n \to 1\}$ n'en est \emph{pas} un~: la suite nulle
        manque (et la somme de deux membres tend vers $2$).
  \item $\{u : u \text{ monotone}\}$ n'en est \emph{pas} un~: $u_n = n$
        et $v_n = -n + (-1)^n$ sont monotones, leur somme
        $(-1)^n$ ne l'est pas~; c'est la stabilité par addition qui
        tombe en défaut, alors même que l'\omterm{def:b1:logic:sets}{ensemble} contient $0$ et
        tous les multiples scalaires de ses membres.
  \item $\{u : u_{n+1} = u_n^2\}$ n'en est \emph{pas} un~: il contient
        $0$ mais $2u$ s'en échappe dès que $u$ en est un membre non
        nul ($2u_{n+1} \neq (2u_n)^2$ en général) --- l'élévation au
        carré est la non-linéarité coupable.
\end{itemize}
L'ordre de travail est toujours le même~: tester $0$ d'abord (le
moins cher), puis la stabilité --- et pour réfuter, un seul
contre-exemple explicite vaut mieux que n'importe quelle quantité de
doute.
\end{example}

\section{Sous-espaces engendrés, sommes, sommes directes}

\begin{definition}[Combinaisons linéaires, sous-espace engendré]\label{def:b1:vspaces:span}
Une \emph{combinaison linéaire} de la famille $(x_1, \dots, x_p)$ de
vecteurs de $E$ est un vecteur $\lambda_1 x_1 + \dots + \lambda_p x_p$
($\lambda_i \in K$). L'\omterm{def:b1:logic:sets}{ensemble} de toutes ces combinaisons est le
\emph{sous-espace engendré}\index{sous-espace engendré}
$\operatorname{Vect}(x_1, \dots, x_p)$~: c'est un \omterm{def:b1:vspaces:subspace}{sous-espace}, le plus
petit contenant la famille.
\end{definition}

\begin{proof}[Démonstration des deux assertions]
Stabilité~: une somme de deux combinaisons linéaires
$\sum\lambda_i x_i + \sum\mu_i x_i = \sum(\lambda_i + \mu_i)x_i$ en
est encore une, et un multiple scalaire
$\mu\sum\lambda_i x_i = \sum(\mu\lambda_i) x_i$ aussi~; la
combinaison nulle montre que $0$ y appartient~: le \omterm{def:b1:vspaces:span}{sous-espace
engendré} est un \omterm{def:b1:vspaces:subspace}{sous-espace}. Minimalité~: soit $H$ un \omterm{def:b1:vspaces:subspace}{sous-espace}
contenant $x_1, \dots, x_p$. Par stabilité par multiplication externe,
chaque $\lambda_i x_i \in H$, et par stabilité par addition leur somme
est dans $H$~: toute combinaison linéaire appartient à $H$, c'est-à-dire
$\operatorname{Vect}(x_1, \dots, x_p) \subseteq H$. Ainsi le
\omterm{def:b1:vspaces:span}{sous-espace engendré} est contenu dans tout \omterm{def:b1:vspaces:subspace}{sous-espace} contenant la
famille~: c'est le plus petit.
\end{proof}

\begin{definition}[Somme, somme directe]\label{def:b1:vspaces:sum}
Pour des \omterm{def:b1:vspaces:subspace}{sous-espaces} $F, G$ de $E$~:
\[
F + G = \{\,u + v : u \in F,\ v \in G\,\}
\]
est un \omterm{def:b1:vspaces:subspace}{sous-espace} (le plus petit contenant $F \cup G$). La somme est
\emph{directe}\index{somme directe}, notée $F \oplus G$, lorsque tout
élément de $F + G$ se décompose de façon \emph{unique} en $u + v$~;
de façon équivalente (voir ci-dessous) lorsque $F \cap G = \{0\}$.
Quand $E = F \oplus G$, les \omterm{def:b1:vspaces:subspace}{sous-espaces} sont
\emph{supplémentaires}\index{sous-espaces supplémentaires} dans $E$.
\end{definition}

\begin{example}[Une somme de deux droites]
\label{ex:b1:vspaces:sumoflines}
Dans $\R^3$, posons $F = \operatorname{Vect}\bigl((1,0,1)\bigr)$ et
$G = \operatorname{Vect}\bigl((0,1,1)\bigr)$. Leur somme est
\[
F + G = \{\,a(1,0,1) + b(0,1,1)\,\}
= \{(a,\ b,\ a + b)\} = \{(x, y, z) : z = x + y\},
\]
le plan passant par l'origine et contenant les deux droites. Il est
strictement plus gros que la réunion $F \cup G$ (la simple croix
formée des deux droites)~: le vecteur $(1, 1, 2) = (1,0,1) + (0,1,1)$
est dans la somme mais sur aucune des deux droites. Et $F \cap G =
\{0\}$ (un vecteur commun impose $a(1,0,1) = b(0,1,1)$, dont les deux
premières \omterm{prop:b1:vspaces:coordinates}{coordonnées} forcent $a = b = 0$)~: la somme est directe, et
$F \oplus G$ est exactement ce plan.
\end{example}

\begin{proposition}\label{prop:b1:vspaces:directsum}
$F + G$ est directe si et seulement si $F \cap G = \{0\}$.
\end{proposition}

\begin{proof}
Si un $w \neq 0$ appartient à $F \cap G$~: $w = w + 0 = 0 + w$ sont
deux décompositions de $w$. Réciproquement, si $u + v = u' + v'$ avec
$u, u' \in F$, $v, v' \in G$, alors $u - u' = v' - v$ appartient à
$F \cap G = \{0\}$~: les décompositions sont uniques.
\end{proof}

\begin{method}[Démontrer que $E = F \oplus G$]
\label{met:b1:vspaces:directsum}
Deux choses à vérifier, chacune avec son ouverture standard.
\begin{enumerate}
  \item \emph{Intersection triviale.} Prendre $x \in F \cap G$,
        écrire les deux conditions d'appartenance, et coincer $x =
        0$. (Ne jamais argumenter sur un dessin~: cf.\ les pièges
        ci-dessous.)
  \item \emph{La somme est tout.} Prendre $x \in E$ quelconque
        et \emph{produire} la décomposition $x = f + g$ ---
        soit en devinant $f$ à partir de la cible ($f$ doit vérifier
        la propriété définissant $F$, ce qui en dicte d'ordinaire la
        formule), soit en résolvant le système linéaire exprimant $x$
        sur des générateurs de $F$ et de $G$.
\end{enumerate}
Quand la formule de décomposition est devinée, l'unicité est
automatique par l'étape 1~; quand seule l'existence pose problème,
c'est à l'étape 2 que le travail se trouve. Les deux exemples
ci-dessous font tourner la méthode~: pour les fonctions paires et
impaires, la formule de $f$ est forcée en évaluant l'identité voulue
en $x$ et en $-x$~; pour les \omterm{def:b1:poly:def}{polynômes} s'annulant en un point, en
évaluant en $a$.
\end{method}

\begin{example}\label{ex:b1:vspaces:evenodd}
Dans $\mathcal{F}(\R, \R)$, les fonctions paires $\mathcal{P}$ et les
fonctions impaires $\mathcal{I}$ sont \omterm{def:b1:vspaces:sum}{supplémentaires}~: toute $f$
s'écrit
\[
f(x) = \underbrace{\frac{f(x) + f(-x)}{2}}_{\text{paire}}
+ \underbrace{\frac{f(x) - f(-x)}{2}}_{\text{impaire}},
\]
et une fonction à la fois paire et impaire est nulle. (Appliqué à
$\exp$, c'est le couple $(\cosh, \sinh)$ du \cref{ch:b1:functions}.)
\end{example}

\begin{example}[{Un couple de supplémentaires dans $K_n[X]$}]
\label{ex:b1:vspaces:hyperplane}
Fixons $a \in K$ et posons $F = \{P \in K_n[X] : P(a) = 0\}$, $G =
\operatorname{Vect}(1)$ (les constantes). Alors $K_n[X] = F \oplus
G$. En effet $F \cap G$ est formé des constantes s'annulant en $a$,
c'est-à-dire $\{0\}$~; et tout $P$ se décompose en
\[
P = \underbrace{\bigl(P - P(a)\bigr)}_{\in F}
+ \underbrace{P(a)}_{\in G} .
\]
La décomposition mérite d'être mémorisée~: retrancher la valeur en un
point est la façon standard de projeter sur les «~fonctions s'annulant
en $a$~». Remarquons que $F$ est un gros \omterm{def:b1:vspaces:subspace}{sous-espace} et $G$ un tout
petit~; un couple de \omterm{def:b1:vspaces:sum}{supplémentaires} n'a aucune raison d'être
équilibré.
\end{example}

\begin{example}[Un supplémentaire n'est jamais unique]
\label{ex:b1:vspaces:manysupplements}
Dans $\R^2$, soit $F = \operatorname{Vect}\bigl((1,0)\bigr)$ (l'axe
des $x$). Les deux \omterm{def:b1:vspaces:subspace}{sous-espaces} $G = \operatorname{Vect}\bigl((0,1)\bigr)$ et
$G' = \operatorname{Vect}\bigl((1,1)\bigr)$ sont \omterm{def:b1:vspaces:sum}{supplémentaires} de
$F$~: chacun ne rencontre $F$ qu'en $0$, et chaque couple a pour somme $\R^2$.
Les décompositions d'un même vecteur diffèrent~:
\[
(2,\ 1.5) = \underbrace{(2, 0)}_{\in F} +
\underbrace{(0, 1.5)}_{\in G}
= \underbrace{(0.5,\ 0)}_{\in F} +
\underbrace{(1.5,\ 1.5)}_{\in G'} .
\]
En fait \emph{toute} droite autre que $F$ elle-même est un
\omterm{def:b1:vspaces:sum}{supplémentaire} de $F$ dans $\R^2$~: les \omterm{def:b1:vspaces:sum}{supplémentaires} foisonnent,
et parler du «~\omterm{def:b1:vspaces:sum}{supplémentaire}~» n'a aucun sens tant qu'une structure
additionnelle (un produit scalaire, \cref{ch:b1:euclid}) n'en
distingue pas un.
\end{example}

\begin{omfigure}
\begin{tikzpicture}[scale=1.5]
  \draw[->, thin] (-0.4,0) -- (2.9,0) node[below] {$x$};
  \draw[->, thin] (0,-0.4) -- (0,1.9) node[left] {$y$};
  \draw[omDef, very thick] (-0.3,0) -- (2.7,0)
    node[below, pos=0.95, font=\small] {$F$};
  \draw[omThm, very thick] (0,-0.3) -- (0,1.7)
    node[left, pos=0.95, font=\small] {$G$};
  \draw[omProp, very thick] (-0.25,-0.25) -- (1.75,1.75)
    node[right, pos=0.97, font=\small] {$G'$};
  \fill (2,1.5) circle (0.045) node[above right, font=\small]
    {$(2, 1.5)$};
  \draw[omThm, dashed] (2,1.5) -- (2,0);
  \draw[omDef, dashed] (2,1.5) -- (0,1.5);
  \draw[omProp, dashed] (2,1.5) -- (0.5,0);
  \fill (2,0) circle (0.035);
  \fill (0.5,0) circle (0.035);
\end{tikzpicture}

{\small Deux décompositions d'un même point de $\R^2$ selon $F$
(l'axe des $x$)~: avec le \omterm{def:b1:vspaces:sum}{supplémentaire} $G$ (descente verticale) et
avec le \omterm{def:b1:vspaces:sum}{supplémentaire} $G'$ (descente oblique). Les composantes sur
$F$ diffèrent~: une projection dépend de la direction de descente.}
\end{omfigure}

\section{Familles libres, familles génératrices, bases}

\begin{definition}\label{def:b1:vspaces:free}
Une famille $(x_1, \dots, x_p)$ de vecteurs de $E$ est~:
\begin{itemize}
  \item \emph{génératrice} (de $E$) lorsque
        $\operatorname{Vect}(x_1,\dots,x_p) = E$~;
  \item \emph{libre}\index{famille libre} (ses vecteurs sont alors
        \emph{linéairement indépendants}) lorsque
        \[
        \lambda_1 x_1 + \dots + \lambda_p x_p = 0
        \implies \lambda_1 = \dots = \lambda_p = 0 ;
        \]
        sinon \emph{liée}~;
  \item une \emph{base}\index{base} lorsqu'elle est libre et
        génératrice.
\end{itemize}
\end{definition}

\begin{proposition}[Coordonnées]\label{prop:b1:vspaces:coordinates}
$(e_1, \dots, e_n)$ est une \omterm{def:b1:vspaces:free}{base} de $E$ si et seulement si tout
$x \in E$ est de façon \emph{unique} une combinaison $x = \lambda_1 e_1
+ \dots + \lambda_n e_n$~; les scalaires $\lambda_i$ sont les
\emph{coordonnées}\index{coordonnées} de $x$ dans la \omterm{def:b1:vspaces:free}{base}.
\end{proposition}

\begin{proof}
Génératrice $=$ existence de la décomposition. Unicité $=$
liberté~: deux décompositions d'un même $x$ diffèrent d'une
combinaison égale à $0$~; la liberté force tous ses coefficients ---
les différences des \omterm{prop:b1:vspaces:coordinates}{coordonnées} --- à être nuls. Réciproquement, une
combinaison nulle non triviale fournit les deux décompositions
$0 = \sum \lambda_i e_i = \sum 0\,e_i$.
\end{proof}

\begin{example}\label{ex:b1:vspaces:bases}
La \emph{\omterm{def:b1:vspaces:free}{base} canonique} de $K^n$~: $e_i = (0, \dots, 1, \dots, 0)$
($1$ en position $i$). Les monômes $(1, X, X^2, \dots, X^n)$~: une
\omterm{def:b1:vspaces:free}{base} de $K_n[X]$ (liberté~: une combinaison nulle est le \omterm{def:b1:poly:def}{polynôme} nul,
donc tous les coefficients sont nuls, \cref{def:b1:poly:def}). Dans
$\C$ vu sur $\R$~: la \omterm{def:b1:vspaces:free}{base} $(1, \iu)$.
\end{example}

\begin{remark}[Les coordonnées sont un travail d'équipe]
\label{rem:b1:vspaces:teameffort}
La première coordonnée de $x$ dans une \omterm{def:b1:vspaces:free}{base} $(e_1, \dots, e_n)$
dépend de \emph{tous} les vecteurs de la \omterm{def:b1:vspaces:free}{base}, pas seulement de $e_1$.
Dans $\R^2$~: le vecteur $(3, 1)$ a pour première coordonnée $3$ dans
la \omterm{def:b1:vspaces:free}{base} canonique, mais $2$ dans la \omterm{def:b1:vspaces:free}{base}
$\bigl((1,0), (1,1)\bigr)$ --- résolvons $(3,1) = a(1,0) + b(1,1)$~:
$b = 1$, $a = 2$. Changer un seul vecteur de \omterm{def:b1:vspaces:free}{base} rebat \emph{toutes}
les \omterm{prop:b1:vspaces:coordinates}{coordonnées}~; le \cref{ch:b1:matrices} mettra ce brassage en
boîte sous le nom de matrice de passage.
\end{remark}

\begin{example}[Tester une base candidate, de bout en bout]
\label{ex:b1:vspaces:testbasis}
La famille $\mathcal{F} = (1 + X,\ 1 + X^2,\ X + X^2)$ est-elle une
\omterm{def:b1:vspaces:free}{base} de $\R_2[X]$~? Notons $u_1, u_2, u_3$ les trois \omterm{def:b1:poly:def}{polynômes}.
\emph{Liberté}~: une combinaison nulle $a\,u_1 + b\,u_2 + c\,u_3 =
0$ donne, coefficient par coefficient,
\[
a + b = 0, \qquad a + c = 0, \qquad b + c = 0 ;
\]
en soustrayant les deux premières, $b = c$, puis la troisième donne
$2b = 0$~: $a = b = c = 0$, la famille est \omterm{def:b1:vspaces:free}{libre}. \emph{Génératrice}~:
plutôt que de résoudre trois systèmes, remarquons la combinaison
symétrique
\[
u_1 + u_2 - u_3 = (1 + X) + (1 + X^2) - (X + X^2) = 2 ,
\]
d'où $1 = \frac12(u_1 + u_2 - u_3)$~; puis
\[
X = u_1 - 1 = \tfrac12\bigl(u_1 - u_2 + u_3\bigr),
\qquad
X^2 = u_2 - 1 = \tfrac12\bigl(-u_1 + u_2 + u_3\bigr).
\]
Les monômes sont dans le \omterm{def:b1:vspaces:span}{sous-espace engendré}, donc tout y est~:
$\mathcal{F}$ est une \omterm{def:b1:vspaces:free}{base}. En prime, en assemblant les trois formules
on obtient les \omterm{prop:b1:vspaces:coordinates}{coordonnées} de tout $P = \alpha + \beta X + \gamma X^2$~:
\[
P = \frac{\alpha + \beta - \gamma}{2}\,u_1
+ \frac{\alpha - \beta + \gamma}{2}\,u_2
+ \frac{-\alpha + \beta + \gamma}{2}\,u_3 .
\]
(Vérification avec $P = X$~: \omterm{prop:b1:vspaces:coordinates}{coordonnées} $\bigl(\frac12,
-\frac12, \frac12\bigr)$, comme trouvé plus haut.) Deux leçons~: une
symétrie dans la famille cache d'ordinaire une combinaison
raccourci~; et une fois la dimension disponible
(\cref{ch:b1:findim}), toute la moitié «~génératrice~» de ce travail
sera gratuite --- trois vecteurs \omterm{def:b1:vspaces:free}{libres} d'un espace de dimension $3$
forment toujours une \omterm{def:b1:vspaces:free}{base}.
\end{example}

\begin{proposition}[Critères de liberté utiles]\label{prop:b1:vspaces:freecriteria}
\begin{enumerate}
  \item Une famille de \emph{\omterm{def:b1:poly:def}{polynômes} non nuls de degrés deux à deux
        distincts} est \omterm{def:b1:vspaces:free}{libre}.
  \item Ajouter un vecteur à une \omterm{def:b1:vspaces:free}{famille libre} la laisse \omterm{def:b1:vspaces:free}{libre} si et
        seulement si ce vecteur est hors du \omterm{def:b1:vspaces:span}{sous-espace engendré} par
        la famille.
  \item Toute sous-famille d'une \omterm{def:b1:vspaces:free}{famille libre} est \omterm{def:b1:vspaces:free}{libre}~; toute
        famille contenant une famille génératrice est génératrice.
\end{enumerate}
\end{proposition}

\begin{proof}
(1) Dans une combinaison nulle, regardons le plus haut degré
présent~: son coefficient doit être nul (rien ne peut annuler ce
degré), puis on descend en cascade.

(2) Si $x \in \operatorname{Vect}(x_1, \dots, x_p)$, la relation $x
- \sum\lambda_i x_i = 0$ est non triviale. Réciproquement, une
combinaison nulle non triviale de $(x_1, \dots, x_p, x)$ doit faire
intervenir $x$ avec un coefficient non nul (sinon elle contredirait la
liberté de la petite famille), et résoudre en $x$ le place dans le
\omterm{def:b1:vspaces:span}{sous-espace engendré}.

(3) Sous-famille~: une combinaison nulle de la sous-famille en est une
de la famille entière, les coefficients manquants étant pris nuls~; la
liberté de la grande famille les tue tous. Sur-famille~: tout vecteur
de $E$ est déjà combinaison de la partie génératrice~; donnons aux
vecteurs \omterm{def:b1:vspaces:sum}{supplémentaires} le coefficient $0$.
\end{proof}

\begin{example}[Le principe de l'escalier]
\label{ex:b1:vspaces:staircase}
Soient $P_0, P_1, \dots, P_n \in K_n[X]$ avec $\deg P_k = k$ pour
chaque $k$ (un «~escalier~» de degrés). Alors $(P_0, \dots, P_n)$ est
une \omterm{def:b1:vspaces:free}{base} de $K_n[X]$. La liberté résulte de la
\cref{prop:b1:vspaces:freecriteria}~(1). Pour le caractère
générateur, raisonnons par descente finie sur le degré~: soit $Q \in
K_n[X]$, $Q \neq 0$, de degré $d$, de coefficient dominant $a$, et
soit $b \neq 0$ le coefficient dominant de $P_d$. Alors $Q -
\frac ab P_d$ est de degré $< d$ (les termes de tête s'annulent)~; en
remplaçant $Q$ par cette différence et en itérant, on atteint après au
plus $n + 1$ étapes le \omterm{def:b1:poly:def}{polynôme} nul, et en remontant les
soustractions on exprime $Q$ comme combinaison des $P_k$. Deux
escaliers déjà rencontrés~: les puissances translatées
$\bigl((X-a)^k\bigr)_{0 \leq k \leq n}$ (\cref{exo:b1:vspaces:4}), et
les produits de Newton
$\bigl((X - x_0)(X - x_1)\cdots(X - x_{k-1})\bigr)_{0 \leq k \leq
n}$, mis à profit dans le devoir maison.
\end{example}

\begin{example}[Liberté dans les espaces de fonctions]\label{ex:b1:vspaces:functionfree}
Dans $\mathcal{F}(\R,\R)$, la famille $(\eu^{a_1 x}, \dots, \eu^{a_p
x})$ avec $a_1 < \dots < a_p$ est \omterm{def:b1:vspaces:free}{libre}~: divisons une combinaison
nulle par $\eu^{a_p x}$ et faisons $x \to +\infty$~; le dernier
coefficient meurt, et l'on descend en cascade
(l'\cref{exo:b1:vspaces:8} détaille cela et ses variantes). La liberté
de fonctions se démontre en \emph{évaluant}~: en des points bien
choisis, à l'infini, ou après dérivation.
\end{example}

\begin{example}[Une relation cachée rétrécit un sous-espace engendré]
\label{ex:b1:vspaces:hiddenrelation}
Dans $\mathcal{F}(\R, \R)$, que vaut
$\operatorname{Vect}(1,\ \cos^2,\ \sin^2)$~? L'identité
$\cos^2 + \sin^2 = 1$ est une combinaison nulle non triviale
\[
1\cdot\mathbf{1} + (-1)\cos^2 + (-1)\sin^2 = 0 :
\]
la famille est liée, et le \omterm{def:b1:vspaces:span}{sous-espace engendré} est déjà engendré par
la seule famille $(1, \cos^2)$ ($\sin^2 = 1 - \cos^2$). Cette famille plus petite
est \omterm{def:b1:vspaces:free}{libre}~: $a + b\cos^2 x = 0$ pour tout $x$ donne, en $x = 0$ et
$x = \frac\pi2$~: $a + b = 0$ et $a = 0$. Le \omterm{def:b1:vspaces:span}{sous-espace engendré} est
donc un \emph{plan} à l'\omterm{def:b1:topology:closure}{intérieur} de l'espace des fonctions --- et il
contient aussi $\cos 2x = 2\cos^2 x - 1$~: les familles de fonctions
trigonométriques d'allure linéaire s'effondrent régulièrement sous
l'effet d'identités, et c'est pourquoi la liberté doit être
\emph{démontrée}, jamais déduite de la longueur de la liste.
\end{example}

\begin{remark}[Pièges classiques]
\label{rem:b1:vspaces:pitfalls}
Quatre pièges classiques.
\emph{Deux à deux ne suffit pas}~: dans $\R^2$, les vecteurs $(1,0)$,
$(0,1)$, $(1,1)$ sont deux à deux non proportionnels, et pourtant liés
--- la liberté est une propriété de la famille \emph{entière},
testée par une unique combinaison globale, jamais deux à deux.
\emph{Le vecteur nul empoisonne tout}~: toute famille contenant
$0$ est liée ($1\cdot 0 = 0$ est une relation non triviale), si
innocents que soient les autres vecteurs.
\emph{Réunion n'est pas somme}~: $F \cup G$ n'est presque jamais un
\omterm{def:b1:vspaces:subspace}{sous-espace} (\cref{def:b1:vspaces:subspace})~; le plus petit
\omterm{def:b1:vspaces:subspace}{sous-espace} contenant les deux est $F + G$, en général bien plus gros
que la réunion --- dans $\R^2$, deux droites distinctes ont pour
réunion une croix, et pour somme le plan tout entier.
\emph{«~Directe~» exige une intersection triviale, pas la
disjonction}~: deux \omterm{def:b1:vspaces:subspace}{sous-espaces} ne sont jamais disjoints (tous deux
contiennent $0$)~; la bonne condition est $F \cap G = \{0\}$, et elle
doit être \emph{démontrée}, pas lue sur un dessin --- cf.\
\cref{ex:b1:vspaces:manysupplements}, où bien des $G$ différents
conviennent.
\emph{La liberté dépend des scalaires}~: le couple $(1, \iu)$ est
\omterm{def:b1:vspaces:free}{libre} dans $\C$ vu comme $\R$-espace vectoriel, mais lié dans $\C$
vu comme $\C$-espace vectoriel ($\iu\cdot 1 + (-1)\cdot\iu = 0$).
Toujours savoir quel \omterm{def:b1:structures:field}{corps} agit avant de déclarer une \omterm{def:b1:vspaces:free}{famille libre}
--- le devoir maison du \cref{ch:b1:findim} transforme exactement
cette sensibilité en démonstrations d'irrationalité.
\end{remark}

\begin{remark}[Où mène ce langage]
\label{rem:b1:vspaces:whereused}
Tout ce qui suit ce chapitre parle le langage installé ici.
Le \cref{ch:b1:findim} compte les vecteurs d'une \omterm{def:b1:vspaces:free}{base} et transforme
«~\omterm{def:b1:vspaces:free}{libre}~» et «~génératrice~» en inégalités sur un unique entier, la
dimension. Le \cref{ch:b1:linmaps} étudie les \omterm{def:b1:logic:map}{applications}
compatibles avec les deux opérations~; les \omterm{def:b1:vspaces:sum}{sommes directes} y
deviennent des projecteurs. Le \cref{ch:b1:matrices} code les vecteurs
par leurs \omterm{prop:b1:vspaces:coordinates}{coordonnées} dans une \omterm{def:b1:vspaces:free}{base} --- la \cref{prop:b1:vspaces:coordinates}
est ce qui autorise ce codage --- et le \cref{ch:b1:euclid}
ajoute longueurs et angles par-dessus la structure linéaire. Dans le
volume de Licence~2, les mêmes axiomes, mot pour mot, valent sur un
\omterm{def:b1:structures:field}{corps} quelconque et en dimension infinie~; rien dans ce chapitre n'a
utilisé la finitude où que ce soit.
\end{remark}

\begin{remark}[Trois fils à suivre dans le livre 3]
\label{rem:b1:vspaces:threads}
Regardez grandir trois idées précises de ce chapitre.
\emph{Le principe de l'escalier}
(\cref{ex:b1:vspaces:staircase}) réapparaît sous la forme de la \omterm{def:b1:vspaces:free}{base}
de Newton dans le devoir maison de ce chapitre, sous celle de la \omterm{def:b1:vspaces:free}{base}
binomiale $(B_k)$ au même endroit, et sous celle de l'astuce de
l'alternant polynomial dans le devoir maison du \cref{ch:b1:det}~: un
lemme, trois dividendes sans le moindre déterminant.
\emph{L'évaluation comme test de liberté}
(\cref{ex:b1:vspaces:functionfree}) devient l'isomorphisme
d'interpolation du \cref{ch:b1:linmaps}, puis le critère de
Vandermonde du \cref{ch:b1:det}, puis le test de Gram du
\cref{ch:b1:euclid}~: un même réflexe, aiguisé trois fois.
\emph{Les \omterm{def:b1:vspaces:sum}{sommes directes}} (\cref{def:b1:vspaces:sum}) deviennent les
projecteurs du \cref{ch:b1:linmaps}, les décompositions orthogonales
$E = F \oplus F^\perp$ du \cref{ch:b1:euclid}, et la décomposition
«~expliqué plus résidu~» des moindres carrés dans le devoir maison du
\cref{ch:b1:multivar}. Bien peu de choses, dans ce livre, ne sont pas,
au fond, l'une de ces trois idées vêtue de neuf.
\end{remark}

\section{Exercices}

\begin{exercise}[$\star$]\label{exo:b1:vspaces:1}
Lesquels des \omterm{def:b1:logic:sets}{ensembles} suivants sont des \omterm{def:b1:vspaces:subspace}{sous-espaces vectoriels}~?
\begin{enumerate}
  \item $\{(x, y, z) \in \R^3 : x + 2y - z = 0\}$~;
  \item $\{(x, y, z) \in \R^3 : x + 2y - z = 1\}$~;
  \item $\{(x, y) \in \R^2 : xy \geq 0\}$~;
  \item $\{P \in \R[X] : P(1) = 0\}$~;
  \item $\{f \in \mathcal{F}(\R,\R) : f \text{ bornée}\}$.
\end{enumerate}
\end{exercise}

\begin{exercise}[$\star$]\label{exo:b1:vspaces:2}
Dans $\R^3$, le vecteur $(1, 2, 1)$ est-il dans
$\operatorname{Vect}\bigl((1,0,1),\, (1,1,0)\bigr)$~? Et $(2, 1,
1)$~? Décrivez $\operatorname{Vect}\bigl((1,0,1),(1,1,0)\bigr)$ par
une équation.
\end{exercise}

\begin{exercise}[$\star$]\label{exo:b1:vspaces:3}
Décidez de la liberté dans $\R^3$~: $\;\bigl((1,1,0), (1,0,1),
(0,1,1)\bigr)$~; $\;\bigl((1,2,3), (2,4,6)\bigr)$~;
$\;\bigl((1,0,0), (1,1,0), (1,1,1), (0,1,1)\bigr)$.
\end{exercise}

\begin{exercise}[$\star$]\label{exo:b1:vspaces:4}
Démontrez que $(1, X - 1, (X-1)^2, (X-1)^3)$ est une \omterm{def:b1:vspaces:free}{base} de
$\R_3[X]$, et donnez les \omterm{prop:b1:vspaces:coordinates}{coordonnées} de $X^3$ dans cette \omterm{def:b1:vspaces:free}{base}.
\emph{(Taylor en $1$~!)}
\end{exercise}

\begin{exercise}[$\star\star$]\label{exo:b1:vspaces:5}
Dans $\R^4$, soient $F = \{(x,y,z,t) : x = y = z\}$ et $G = \{(x,y,z,t)
: x = t = 0\}$. Démontrez que $F \oplus G = \R^4$, et décomposez
$(1,2,3,4)$ en conséquence.
\end{exercise}

\begin{exercise}[$\star\star$]\label{exo:b1:vspaces:6}
Dans l'espace des suites, soit $F$ l'\omterm{def:b1:logic:sets}{ensemble} des suites convergentes
et $G = \operatorname{Vect}(u)$ où $u_n = (-1)^n$.
Démontrez que $F \cap G = \{0\}$. La somme $F + G$ est-elle l'espace
des suites tout entier~?
\end{exercise}

\begin{exercise}[$\star\star$]\label{exo:b1:vspaces:7}
Soient $F, G, H$ des \omterm{def:b1:vspaces:subspace}{sous-espaces} de $E$. Démontrez que
\[
F \cap (G + (F \cap H)) = (F \cap G) + (F \cap H),
\]
et montrez sur un exemple dans $\R^2$ que la distributivité sans
restriction $F \cap (G + H) = (F\cap G) + (F \cap H)$ est fausse.
\end{exercise}

\begin{exercise}[$\star\star\star$]\label{exo:b1:vspaces:8}
Démontrez que les familles suivantes de $\mathcal{F}(\R, \R)$ sont
\omterm{def:b1:vspaces:free}{libres}~:
\begin{enumerate}
  \item $(\eu^{a_1 x}, \dots, \eu^{a_p x})$ pour $a_1 < \dots < a_p$~;
  \item $(\cos x, \sin x, \cos 2x, \sin 2x)$~;
  \item $(x \mapsto \abs{x - a_1}, \dots, x \mapsto \abs{x - a_p})$
        pour des $a_i$ distincts \emph{(la \omterm{def:b1:derivative:def}{dérivabilité} tombe en
        défaut en exactement un point par fonction)}.
\end{enumerate}
\end{exercise}

\begin{exercise}[$\star\star\star$]\label{exo:b1:vspaces:9}
Soient $E$ un $K$-espace vectoriel et $F, G, H$ des \omterm{def:b1:vspaces:subspace}{sous-espaces} avec
$F + G = F + H$, $F \cap G = F \cap H$ et $G \subseteq H$. Démontrez
que $G = H$. Donnez un contre-exemple sans l'hypothèse
$G \subseteq H$.
\end{exercise}

\begin{exercise}[$\star\star$]\label{exo:b1:vspaces:10}
Dans $\R[X]$, soit $\mathcal P$ l'\omterm{def:b1:logic:sets}{ensemble} des \omterm{def:b1:poly:def}{polynômes} \emph{pairs}
($P(-X) = P(X)$) et $\mathcal I$ celui des \omterm{def:b1:poly:def}{polynômes} \emph{impairs}
($P(-X) = -P(X)$). Démontrez que $\R[X] = \mathcal P \oplus \mathcal
I$, et montrez que $\mathcal P =
\operatorname{Vect}(1, X^2, X^4, \dots)$, c'est-à-dire que les
\omterm{def:b1:poly:def}{polynômes} pairs sont exactement les \omterm{def:b1:poly:def}{polynômes} en $X^2$.
\end{exercise}

\begin{exercise}[$\star\star$]\label{exo:b1:vspaces:11}
Soit $(x_1, x_2, x_3)$ une \omterm{def:b1:vspaces:free}{famille libre} d'un \omterm{def:b1:vspaces:def}{espace vectoriel} réel $E$.
Démontrez que $(x_1 + x_2,\ x_2 + x_3,\ x_3 + x_1)$ est \omterm{def:b1:vspaces:free}{libre}. La
famille analogue de quatre vecteurs $(x_1 + x_2,\ x_2 + x_3,\ x_3 +
x_4,\ x_4 + x_1)$ est-elle \omterm{def:b1:vspaces:free}{libre} lorsque $(x_1, x_2, x_3, x_4)$ l'est~?
\end{exercise}

\begin{exercise}[$\star\star\star$]\label{exo:b1:vspaces:12}
Soient $E$ un \omterm{def:b1:vspaces:def}{espace vectoriel} sur $\R$ (ou $\C$) et $F_1, \dots,
F_k$ des \omterm{def:b1:vspaces:subspace}{sous-espaces} \emph{propres} de $E$ (chaque $F_i \neq E$).
\begin{enumerate}
  \item Traitez directement le cas $k = 2$~: si $F_1 \not\subseteq
        F_2$ et $F_2 \not\subseteq F_1$, prenez $x \in F_1
        \setminus F_2$ et $y \in F_2 \setminus F_1$ et localisez $x
        + y$.
  \item Démontrez en général que $E \neq F_1 \cup \dots \cup F_k$~: un
        \omterm{def:b1:vspaces:def}{espace vectoriel} sur un \omterm{def:b1:structures:field}{corps} infini n'est jamais réunion
        finie de \omterm{def:b1:vspaces:subspace}{sous-espaces} propres. \emph{(Prenez $k$ minimal,
        choisissez $x \in F_1$ hors des autres $F_i$, choisissez $y
        \notin F_1$, et suivez la droite $t \mapsto y + tx$.)}
\end{enumerate}
\end{exercise}

\section{Problème~: interpolation, trois bases pour un même espace}

\begin{problem}\label{pb:b1:vspaces:1}
Fixons $n + 1$ points \emph{distincts} $x_0, x_1, \dots, x_n$ de $\R$.
Ce problème revisite l'\omterm{thm:b1:poly:lagrange}{interpolation de Lagrange}
(\cref{thm:b1:poly:lagrange}) avec les yeux de ce chapitre~: l'espace
$\R_n[X]$ porte trois \omterm{def:b1:vspaces:free}{bases} naturelles --- celle de Lagrange, celle de
Newton, et (pour des points équidistants) la \omterm{def:b1:vspaces:free}{base} binomiale ---
et chacune rend une question facile. La route s'achève sur un
véritable théorème d'arithmétique~: la caractérisation, due à Pólya,
des \omterm{def:b1:poly:def}{polynômes} envoyant $\Z$ dans $\Z$.\index{interpolation de Lagrange}
\index{différences divisées}\index{polynôme à valeurs entières}

\medskip
\noindent\textbf{Partie I --- La \omterm{def:b1:vspaces:free}{base} de Lagrange.} Pour $0 \leq i
\leq n$, posons
\[
L_i \;=\; \prod_{j \neq i} \frac{X - x_j}{x_i - x_j} \;\in\;
\R_n[X].
\]
\begin{enumerate}
  \item Vérifiez que $\deg L_i = n$ et que $L_i(x_j) = 1$ si $j =
        i$, et $0$ si $j \neq i$.
  \item Démontrez que la famille $(L_0, \dots, L_n)$ est \omterm{def:b1:vspaces:free}{libre}.
  \item Démontrez que pour tout $P \in \R_n[X]$,
        \[
        P \;=\; \sum_{i=0}^{n} P(x_i)\, L_i ,
        \]
        et concluez que $(L_0, \dots, L_n)$ est une \omterm{def:b1:vspaces:free}{base} de
        $\R_n[X]$. \emph{(Considérez la différence des deux membres
        et comptez ses racines, \cref{cor:b1:poly:nroots}.)}
  \item Déduisez-en le théorème d'interpolation~: pour toutes valeurs
        $y_0, \dots, y_n \in \R$, il existe un \emph{unique} $P \in
        \R_n[X]$ tel que $P(x_i) = y_i$ pour tout $i$. Dans la \omterm{def:b1:vspaces:free}{base}
        de Lagrange, quelles sont les \omterm{prop:b1:vspaces:coordinates}{coordonnées} d'un \omterm{def:b1:poly:def}{polynôme} $P$~?
  \item Démontrez les identités
        \[
        \sum_{i=0}^{n} L_i = 1
        \qquad\text{et, pour } 0 \leq k \leq n,\qquad
        \sum_{i=0}^{n} x_i^{k}\, L_i = X^{k} .
        \]
\end{enumerate}

\medskip
\noindent\textbf{Partie II --- La \omterm{def:b1:vspaces:free}{base} de Newton et les \omterm{pb:b1:vspaces:1}{différences
divisées}.} Posons $N_0 = 1$ et $N_k = (X - x_0)(X - x_1) \cdots
(X - x_{k-1})$ pour $1 \leq k \leq n$. Pour une fonction $f$ définie
aux nœuds, définissons les \emph{\omterm{pb:b1:vspaces:1}{différences divisées}} par
$f[x_i] = f(x_i)$ et
\[
f[x_i, \dots, x_{i+k}] \;=\;
\frac{f[x_{i+1}, \dots, x_{i+k}] - f[x_i, \dots, x_{i+k-1}]}
{x_{i+k} - x_i} .
\]
\begin{enumerate}[resume]
  \item Démontrez que $(N_0, N_1, \dots, N_n)$ est une \omterm{def:b1:vspaces:free}{base} de
        $\R_n[X]$.
  \item Calculez $f[x_0, x_1]$ et $f[x_0, x_1, x_2]$ en fonction des
        valeurs de $f$, puis calculez toutes les \omterm{pb:b1:vspaces:1}{différences divisées}
        de $f(x) = x^2$ en trois nœuds quelconques.
  \item (Lemme d'Aitken) Soient $R$ interpolant $f$ en $x_0, \dots,
        x_{n-1}$ et $Q$ interpolant $f$ en $x_1, \dots, x_n$,
        tous deux de degré $\leq n - 1$. Démontrez que
        \[
        S \;=\; \frac{(X - x_0)\,Q - (X - x_n)\,R}{x_n - x_0}
        \]
        interpole $f$ en $x_0, x_1, \dots, x_n$.
  \item Déduisez-en, par récurrence sur le nombre de nœuds, que le
        coefficient de $X^{k}$ dans l'interpolant de $f$ en $x_0,
        \dots, x_k$ vaut exactement $f[x_0, \dots, x_k]$.
  \item Démontrez la \emph{formule d'interpolation de Newton}~:
        l'interpolant de $f$ en $x_0, \dots, x_n$ est
        \[
        P \;=\; \sum_{k=0}^{n} f[x_0, \dots, x_k]\, N_k ,
        \]
        et déduisez-en la formule close
        \[
        f[x_0, \dots, x_k] \;=\; \sum_{i=0}^{k}
        \frac{f(x_i)}{\prod_{j \neq i,\, j \leq k} (x_i - x_j)} ,
        \]
        qui montre que $f[x_0, \dots, x_k]$ ne dépend pas de l'ordre
        des nœuds.
\end{enumerate}

\medskip
\noindent\textbf{Partie III --- Nœuds équidistants~: l'opérateur
de différence.} Désormais les nœuds sont $0, 1, 2, \dots$ et, pour
un \omterm{def:b1:poly:def}{polynôme} $P$, on pose
\[
\Delta P(X) = P(X + 1) - P(X),
\qquad
B_k = \frac{X(X-1)\cdots(X-k+1)}{k!} \quad (B_0 = 1).
\]
\begin{enumerate}[resume]
  \item Montrez que si $\deg P = m \geq 1$ de coefficient dominant
        $a$, alors $\deg \Delta P = m - 1$ de coefficient dominant
        $m\,a$, et que $\Delta$ tue les constantes.
  \item Montrez que $(B_0, B_1, \dots, B_n)$ est une \omterm{def:b1:vspaces:free}{base} de
        $\R_n[X]$ et que $\Delta B_k = B_{k-1}$ pour $k \geq 1$.
  \item (Formule des différences progressives de Newton) Démontrez que
        tout $P \in \R_n[X]$ vérifie
        \[
        P \;=\; \sum_{k=0}^{n} \bigl(\Delta^{k} P\bigr)(0)\, B_k .
        \]
  \item Démontrez que pour tout $k \geq 0$,
        \[
        \bigl(\Delta^{k} P\bigr)(0) \;=\;
        \sum_{j=0}^{k} (-1)^{k-j} \binom{k}{j} P(j) .
        \]
  \item Montrez que si $\deg P = n$ de coefficient dominant $a_n$,
        alors $\Delta^{n} P$ est la constante $n!\,a_n$ et
        $\Delta^{n+1} P = 0$.
\end{enumerate}

\medskip
\noindent\textbf{Partie IV --- \omterm{pb:b1:vspaces:1}{Polynômes à valeurs entières}.} Un
\omterm{def:b1:poly:def}{polynôme} $P \in \R[X]$ est \emph{à valeurs entières} lorsque $P(m) \in
\Z$ pour tout $m \in \Z$.
\begin{enumerate}[resume]
  \item Démontrez que chaque $B_k$ est à valeurs entières.
        \emph{(Traitez séparément $m \geq k$, $0 \leq m < k$ et $m <
        0$~; pour $m = -q < 0$, montrez que $B_k(-q) = (-1)^k
        \binom{q + k - 1}{k}$.)}
  \item Démontrez la \emph{caractérisation de Pólya}~: $P \in \R_n[X]$
        est à valeurs entières si et seulement si ses \omterm{prop:b1:vspaces:coordinates}{coordonnées} dans
        la \omterm{def:b1:vspaces:free}{base} $(B_0, \dots, B_n)$ sont des entiers.
  \item Déduisez-en~: si $P \in \R_n[X]$ prend des valeurs entières en
        $n + 1$ entiers \emph{consécutifs} $a, a+1, \dots, a+n$, alors
        $P$ est à valeurs entières. \emph{(Translation~: appliquez
        l'étude à $Q(X) = P(X + a)$.)}
  \item Déduisez de la question 16 qu'un produit de $k$ entiers
        consécutifs est toujours divisible par $k!$.
  \item Soit $P = \dfrac{X(X+1)(2X+1)}{6}$. Calculez sa table de
        Newton en $0, 1, 2, 3$, écrivez $P$ dans la \omterm{def:b1:vspaces:free}{base} $(B_k)$, et
        concluez que $P$ est à valeurs entières bien qu'aucun de ses
        coefficients monomiaux ne soit entier. Vérifiez que $\Delta P =
        (X+1)^2$ et déduisez-en que $P(m) = 1^2 + 2^2 + \dots + m^2$
        pour $m \in \N$.
\end{enumerate}

\medskip
\noindent\textbf{Partie V --- Dividendes.}
\begin{enumerate}[resume]
  \item Soit $P \in \R_n[X]$ interpolant les valeurs $2^i$ en $i =
        0, 1, \dots, n$. Montrez que $P = B_0 + B_1 + \dots + B_n$
        et que $P(n + 1) = 2^{n+1} - 1$~: le «~schéma de
        doublement~» se brise toujours au point suivant.
  \item (Primitive discrète) Démontrez que pour tous entiers $m
        \geq 1$ et $k \geq 0$,
        \[
        \sum_{j=0}^{m-1} B_k(j) \;=\; B_{k+1}(m),
        \]
        c'est-à-dire l'identité de la crosse de hockey $\sum_{j=k}^{m-1}
        \binom{j}{k} = \binom{m}{k+1}$.
  \item Développez $X^2$ et $X^3$ dans la \omterm{def:b1:vspaces:free}{base} $(B_k)$ et déduisez-en
        des formules closes pour $\sum_{j=0}^{m-1} j^2$ et
        $\sum_{j=0}^{m-1} j^3$~; retrouvez l'identité de Nicomaque
        $1^3 + \dots + m^3 = (1 + \dots + m)^2$.
  \item Prenez $n = 2$ et les nœuds $0, 1, 2$. Écrivez les
        \omterm{prop:b1:vspaces:coordinates}{coordonnées} de $X^2$ dans les trois \omterm{def:b1:vspaces:free}{bases} de ce problème~: la
        \omterm{def:b1:vspaces:free}{base} monomiale, la \omterm{def:b1:vspaces:free}{base} de Lagrange, la \omterm{def:b1:vspaces:free}{base} de Newton.
        Confrontez les trois réponses aux questions 4 et 9.
  \item Synthèse. En quatre phrases~: quel concept d'\omterm{def:b1:vspaces:def}{espace vectoriel}
        rend la question 4 automatique~; pourquoi la \omterm{def:b1:vspaces:free}{base} de Newton
        calcule les \omterm{prop:b1:vspaces:coordinates}{coordonnées} \emph{de proche en proche} alors que
        la \omterm{def:b1:vspaces:free}{base} de Lagrange les lit \emph{instantanément}~; quel
        critère de liberté les deux \omterm{def:b1:vspaces:free}{bases} partagent~; et en quel sens
        précis le théorème de Pólya dit que le fait qu'un \omterm{def:b1:poly:def}{polynôme}
        soit à valeurs entières est une propriété de ses \omterm{prop:b1:vspaces:coordinates}{coordonnées}
        \emph{dans la bonne \omterm{def:b1:vspaces:free}{base}}.
\end{enumerate}
\end{problem}
